% !TEX root = ../main.tex
\clearpage
\chapter{부록}
\label{ch:appendix}

% Demote headings one level under Chapter 8 (sections become 8.1.x).
\renewcommand{\dissertationSubheading}[2][]{%
  \phantomsection
  \subsection{#2}%
  \if\relax\detokenize{#1}\relax\else\label{#1}\fi
}
\renewcommand{\dissertationSubsubheading}[2][]{%
  \phantomsection
  \subsubsection{#2}%
  \if\relax\detokenize{#1}\relax\else\label{#1}\fi
}

\section{평가 처리 과정과 재현성}
\label{ch:appendix-eval}

이 부록은 \ref{ch:results}장을 뒷받침하는 평가 처리 과정(파이프라인)을 기록해요. 처리 단계, 자료를 꺼내는 추출 설계, 기본 매개변수, 대리 지표(대신 재는 값)의 공식, 그리고 용어 풀이를 담았어요.

\dissertationSubheading[sec:pipeline-stages]{처리 단계}

평가는 네 단계로 이루어져요. 앞의 세 단계는 아비트럼 원(Arbitrum One)에서 관찰 기간 2025-12-01부터 2026-05-31까지(UTC, 양 끝 날짜 포함)의 자료로 돌아가요. 네 번째 단계는 여기에 2026년 6월부터 8월까지의 달을 더해요.

\begin{enumerate}
\item 평판 그래프와 점수: 허락이나 송금 기록을 꺼내기 전에, 먼저 GMX~V2 자료로 매칭 코호트(모든 점수가 똑같이 점수를 매기는 지갑 묶음, $n=5{,}521$)를 정해요. 이 코호트의 \texttt{Approval} 로그와 \texttt{Transfer} 로그(계약이 남기는 기록)를 한 달씩 나누어 가져와요. 질의마다 먼저 미리 돌려 보기를 하고, $150$\,GiB를 넘으면 멈추는 한도를 둬요. 마지막 허락 그래프와 송금 그래프를 만들고, 각각 엔도스랭크(EndorseRank)와 AWP(적응형 가중 페이지랭크)로 점수를 매겨요. 이어서 금액을 그대로 쓴 두 층을 만들고, C-PR(결합 페이지랭크, $\lambda\in\{0,0.25,0.5,0.75,1\}$)과 S-PR(보증 씨앗 페이지랭크) 점수를 매겨요. 모두 \ref{sec:notation}절의 풀이 프로그램(solver) 설정을 써요.
\item 점검, 일치도(점수와 지표의 순서가 얼마나 맞는지), 벤치마크(속도 재기): \ref{sec:proxy-formulas}절의 점검 값을 계산해요. 모든 점수에 대해 스피어먼의 로($\rho$)와 켄달의 타우($\tau$, 두 줄 세우기가 얼마나 같은 순서인지 나타내는 수)를 구하고, 95\% 짝지은 부트스트랩(다시 뽑기) 구간(다시 뽑은 표본 400개, 시드 42)과 함께 보고해요. \ref{sec:statistical-tests}절의 대비(두 계수의 짝지은 차이)도 함께 보고해요. \ref{sec:computational-benchmark}절의 절차대로, 방법마다 풀이 프로그램을 한 번 부르는 데 걸리는 시간을 재요. 그리고 감쇠 계수, 상위 토큰, 표본 정의를 차례로 바꿔 보는 실행도 돌려요.
\item 시간 홀드아웃(나중 기간을 떼어 두고 맞혀 보는 시험): 스펜더(spender, 허락을 받아 대신 쓰는 쪽) 코호트($n=1{,}335$)에서 점수를 2026년 2월 28일에 고정해요. 그리고 2026년 3월부터 5월까지 나타난 새 허락 쌍(새로 생긴 주인–스펜더 허락)과 새 송금자(새로 송금을 보내온 주소)를 세요. 모든 점수와 두 원시 차수(화살표를 그냥 센 수) 기준선(비교 기준)을 두 라벨(나중에 채점에 쓰는 정답 값)에 견주어, 짝지은 다시 뽑기 400번으로 보고해요. 결합 점수를 돌리기 전에 정해 둔 대비와 그 뒤에 계산한 대비도 함께 보고해요(\texttt{scripts/run\_holdout\_eval.py --cohort spenders}).
\item 등록 재현(계획을 먼저 공개해 두고 새 자료로 다시 해 보기): \texttt{config/fresh\_holdout\_2026q3.yaml}(설정 파일)과 그 계획서 \texttt{docs/fresh\_holdout\_2026q3\_plan.md}를 커밋(저장소에 기록)해 둔 상태에서, 먼저 봄 홀드아웃으로 평가 코드를 점검해요(\texttt{scripts/run\_fresh\_holdout.py --spring-check}). 그다음 같은 SQL과 같은 지갑 거르기 조건으로 6월부터 8월까지의 로그를 꺼내요(\texttt{scripts/extract\_fresh\_window.py --dry-run}, 그다음 \texttt{--extract}). 그리고 2026년 5월 31일에 고정한 점수로 두 코호트를 평가해요(\texttt{scripts/run\_fresh\_holdout.py}). 추출 스크립트는 상태를 ‘등록됨’으로 바꾼 등록이 커밋되기 전에는 돌아가지 않아요. 두 스크립트 모두 두 등록 파일의 SHA-256 해시(파일 내용으로 만든 지문 같은 값)를 기록해요.
\end{enumerate}

결합 점수를 하나라도 계산하기 전에, 섞는 비율, 홀드아웃 라벨, 그리고 RQ5(다섯 번째 연구 질문)에 대한 첫 번째 규칙을 설정 파일과 2026년 9월 14일의 계획서에 적어 두었어요. 2026년 9월 30일에는 역위험 대리 지표(위험이 작을수록 값이 커야 하는 지표) 두 개의 부호가 거꾸로 되어 있는 것을 찾아내 바로잡았어요(부록~\ref{sec:supplementary-summary}). 같은 날 \texttt{scripts/classify\_spenders.py}는 점수를 매긴 기간에 어떤 스펜더가 로그를 남겼거나 거래를 보냈는지를 빅쿼리(BigQuery)에서 확인했어요. 그리고 \texttt{scripts/check\_spender\_code.py}는 아비트럼 노드(블록체인 기록을 가진 컴퓨터)에서 스펜더마다 코드를 읽었어요(\texttt{data/2-processed-tables-and-evaluations/spring-holdout-2026-03-to-2026-05/spender\_account\_types.csv}). 2026년 10월 1일에는 등록에 민감도 분석(조건을 바꿔 결과가 얼마나 움직이는지 보기) 두 개가 더해졌어요(\ref{sub:fresh-holdout}절). AWP의 값 매개변수는 설정 파일의 \texttt{reputation.awp\_paper}에 있어요. 그날 늦게 상태를 ‘등록됨’으로 바꾸었어요(\texttt{phd\_works}의 커밋 \texttt{4085ab9}). 이어서 6월부터 8월까지의 로그를 한 달 단위 질의 아홉 개로 꺼냈는데, 요금은 약 1.04\,TB만큼 매겨졌어요. 그리고 \texttt{scripts/run\_fresh\_holdout.py}가 그 로그를 평가했어요. \texttt{data/3-published-results-for-thesis/registered-replication-2026-06-to-2026-08/}에는 추출 목록 파일(매니페스트)과, 표~\ref{tab:fresh-holdout}, \ref{tab:fresh-contrasts}와~\ref{tab:fresh-traders}의 바탕이 된 요약이 들어 있어요.

\ref{ch:results}장의 나머지 표들은 2026년 10월 2일에 클라우드 컨테이너(인터넷 너머에서 빌려 쓰는 컴퓨터 환경)에서 돌린 실행에서 나왔어요. 이 컨테이너에는 2.8\,GHz 인텔 제온(Intel Xeon) 프로세서의 코어 네 개와 15\,GB 메모리가 있었고, Python~3.11과 NumPy~2.4, SciPy~1.17, pandas~3.0을 썼어요. 돌린 실행은 다음과 같아요. \texttt{scripts/run\_dissertation\_eval.py --real --robustness --benchmark-repeats 5}; \texttt{--cohort spenders}와 \texttt{--cohort matched}를 각각 붙인 \texttt{scripts/run\_holdout\_eval.py}; 등록 재현의 코호트와 라벨로 엔도스랭크 점수를 매기는 \texttt{scripts/run\_fresh\_posthoc.py}(표~\ref{tab:fresh-posthoc}); 그리고 표~\ref{tab:tie-sensitivity}와~\ref{tab:holdout-receiving}의 바탕이 된 \texttt{scripts/run\_supplementary\_checks.py}예요. 같은 날 \texttt{scripts/run\_fresh\_holdout.py --spring-check}를 봄 기간에 다시 돌렸더니, 10월 1일에 냈던 값이 그대로 다시 나왔어요. 표는 \texttt{scripts/export\_latex\_results.py}가 \texttt{results/tables/}에 썼고, 그림은 \texttt{scripts/generate\_dissertation\_figures.py}가 만들었어요. 요약 파일들은 그때 쓰던 설정 파일의 사본과 함께 \texttt{data/3-published-results-for-thesis/} 아래에 커밋되어 있어요. 봄 홀드아웃 폴더에는 \texttt{eval\_summary.json}, \texttt{supplementary\_checks.json}, \texttt{eval\_summary\_spenders.json}, \texttt{eval\_summary\_matched\_traders.json}이 있고, 등록 재현 폴더에는 \texttt{fresh\_holdout\_summary.json}과 \texttt{posthoc\_endorserank.json}이 있어요. 그래서 처리 과정을 다시 돌리지 않아도 표의 어떤 숫자든 그 요약과 맞춰 볼 수 있어요. 공개 저장소 \url{https://github.com/abitocodes/phd_works}에는 \texttt{[revision-pending-commit]} 버전을 기준으로 다음이 들어 있어요. \texttt{2-Dissertation-Draft-Works/wallet-reputation-experiments}에 있는 평가 소스 코드(이 폴더의 이름은 2026년 10월 1일까지 \texttt{A-Skill-Programs/margin\_rank}였고, 커밋 \texttt{4085ab9}에서도 그 이름이에요), 버전 관리되는 설정 파일(\texttt{config/margin\_config.yaml}), 꺼낸 자료, 그리고 요약 파일들이에요. 그 폴더의 \texttt{data/README.md} 파일이 자료 파일 하나하나를 설명해요.

실제로 자료를 꺼내기 전에, 구현한 코드는 오프라인(인터넷 연결 없이)으로 점검해요. 합성 시험(가짜 자료로 하는 시험)은 클라우드에 접속하지 않고, 시험용으로 미리 만든 그래프 위에서 같은 점수 매기기 논리와 일치도 계산 논리를 돌려요. 실제 추출에는 구글 클라우드(Google Cloud) 인증 정보와 요금을 낼 프로젝트가 필요해요. 그리고 읽는 범위는 공개된 아비트럼 원 로그 데이터셋으로만 한정돼요.

\dissertationSubheading[sec:bq-extraction]{빅쿼리 추출 설계}

ERC-20(토큰 공통 규칙)의 \texttt{Approval} 이벤트와 \texttt{Transfer} 이벤트는 저마다 표준 topic-$0$ 해시(이벤트 종류를 알려 주는 첫 번째 꼬리표 값)로 알아봐요. 질의는 주인/스펜더 주소나 보낸 쪽/받는 쪽 주소가 평가 지갑 목록에 있는 로그만 읽도록 제한해요. 그래서 읽는 비용이 체인 전체가 아니라 코호트 크기에 묶여요. 한 달씩 나누어 질의하면, 모든 질의가 운영상 정한 바이트 상한 아래에 머물러요. 부록~\ref{ch:appendix-sql}에서 대표적인 SQL 모양을 보여 줘요.

나중에 검사(감사)할 수 있도록, 이 연구의 수집 기록에는 추출 기록(처리한 바이트 수, 행 수, 시각)이 남아 있어요. 이 연구는 그 결과로 나온 코호트 크기와 화살표(엣지) 수만 보고하고, 운영 기록은 보고하지 않아요.

확장 지갑 풀(더 큰 주소 모음, $N=99{,}080$; \ref{sub:expanded-pool}절)에 더해 넣은 주소 $93{,}559$개에 대해서는 두 번째 단계의 추출을 계획했었어요. 이 주소들에 대한 허락·송금 월별 질의 여섯 개는 미리 돌려 보기 추정치가 $1.3$\,TiB를 넘었어요. 그리고 결과를 저장 공간을 거쳐 내려받는 단계가 여러 번 실패했어요. 그래서 연구 기간 안에 이 추출을 그만두었어요. 그 결과, 더해 넣은 주소들은 확장 풀 실험에 뽑힌 주소로만 들어가요. 이 주소들에는 자기 허락 화살표나 송금 화살표가 없어요. 또 풀이 프로그램의 그래프는 남겨 둔 화살표의 양 끝에서 점(노드)을 가져오므로, 이 주소들은 점도 더하지 않아요. 그래서 표~\ref{tab:benchmark-scaling}의 전체 풀 행은 매칭 코호트 그래프를 그대로 다시 만들어 내요. \ref{ch:conclusion}장은 이 추출을 마치는 일을 앞으로 할 일로 적어 두었어요.

\dissertationSubheading[sec:scaling-appendix]{규모에 따른 변화 표}

표~\ref{tab:benchmark-scaling-incohort}는 코호트 안에서 $n \leq 5{,}521$일 때 실행 시간이 규모에 따라 어떻게 바뀌는지 보여 줘요. 시간은 \ref{sec:computational-benchmark}절의 시간 재기 절차로 쟀어요. 이 표의 마지막 행은 표~\ref{tab:benchmark-runtime}에 있는 것과 같은 측정이에요. 코호트 안의 단계는 모두 화살표를 꺼낸 지갑들의 부분 집합이에요. 그래서 $|V|$(점의 수)와 $|E|$(화살표의 수)가 둘 다 $n$과 함께 커져요. 확장 풀 단계에서도 둘 다 커지지만, 매칭 코호트 그래프의 크기까지만 커져요.

% Real BigQuery data -- regenerate with:
%   2-Dissertation-Draft-Works/wallet-reputation-experiments/scripts/run_dissertation_eval.py --real --export-latex
% Generated by export_latex_results.py -- do not edit by hand unless re-export fails
\begin{table}[htbp]
\centering
\caption[In-cohort runtime scaling.]{In-cohort runtime scaling (SHA256 seed \texttt{benchmark-scale-v1}; $n \leq 5{,}521$ matched GMX cohort). Here both $|V|$ and $|E|$ grow with $n$ because every stage is a subset of wallets whose edges were extracted. Runtime = mean wall-clock seconds over 5 timed PageRank solves after one untimed warm-up; SD over the same 5 runs.}
\label{tab:benchmark-scaling-incohort}
\fitwidth{%
\begin{tabular}{rrrrrrrr}
\toprule
$n$ & ER (s) & AWP (s) & Speedup & ER $|V|$ & ER $|E|$ & AWP $|V|$ & AWP $|E|$ \\
\midrule
1{,}000 & 0.096 & 0.821 & 8.6$\times$ & 1{,}302 & 2{,}584 & 6{,}463 & 23{,}058 \\
2{,}000 & 0.185 & 1.755 & 9.5$\times$ & 2{,}362 & 5{,}147 & 11{,}849 & 47{,}031 \\
4{,}000 & 0.392 & 3.645 & 9.3$\times$ & 4{,}645 & 10{,}557 & 22{,}591 & 97{,}105 \\
5{,}521 & 0.547 & 5.029 & 9.2$\times$ & 6{,}442 & 14{,}727 & 30{,}791 & 137{,}087 \\
\bottomrule
\end{tabular}
}
\end{table}


\dissertationSubheading[sec:delegation-deferred]{신용 위임 추출(미룸)}

Aave~V3(돈을 빌려주고 빌리는 디파이 서비스)는 신용 위임(내가 빌릴 수 있는 권한을 남에게 맡기는 일)을 \texttt{BorrowAllowanceDelegated} 이벤트로 기록해요. 이것은 대신 쓸 수 있는 권한을 맡기는 두 번째 온체인(블록체인 위의) 방식이에요. 이 연구는 이 이벤트를 꺼내지 않았어요. 그래서 보관해 둔 행렬(부록~\ref{ch:appendix-archive})에서 위임에 바탕을 둔 행에는 화살표가 없어요. 이 이벤트를 허락 그래프에 더하는 일은 앞으로 할 일이에요.

\dissertationSubheading[sec:supplementary-summary]{보충 일치도 요약 표}

처리 과정은 지표 가족(비슷한 지표 묶음) 여섯 개를 계산해요. 표~\ref{tab:alignment-summary}는 모든 점수의 가족 평균 $\tau$를 가족마다 보여 줘요. 그 가운데에는 청산(손실이 너무 커져 거래가 강제로 끝나는 일), 역위험, GMX 거래 성공 가족이 있어요. 이 세 가족은 GF-PR 참고 점수와 함께 주된 분석에서 뺐어요(\ref{sub:gf-pr}절). 2026년 9월 30일까지 처리 과정은 역위험 대리 지표 세 개 가운데 두 개, 곧 손실 합계와 가장 나쁜 닫기(손실이 가장 컸던 포지션 닫기)의 부호를 뒤집었어요. 그래서 두 지표 모두 손실이 가장 큰 지갑을 가장 높게 줄 세웠어요. 지금은 두 지표도 다른 모든 대리 지표와 같은 방향으로 맞추어서, 0에 가까운 값일수록 손실이 작다는 뜻이에요. 표에는 바로잡은 값이 나와 있고, 이렇게 바로잡아도 다른 칸은 하나도 바뀌지 않아요. 이 가족에서는 모든 점수가 음수예요. 이것은 손실 합계와 가장 나쁜 닫기가 지갑이 거래를 많이 할수록 커진다는 점과 맞아요(\ref{sec:same-window-alignment}절). 송금, 허락, 시빌(가짜 주소를 많이 만들어 속이는 공격) 안정성 열은 표~\ref{tab:alignment-hybrid}의 값을 구간 없이 다시 보여 줘요.

% Real BigQuery data -- regenerate with:
%   2-Dissertation-Draft-Works/wallet-reputation-experiments/scripts/run_dissertation_eval.py --real --export-latex
% Generated by export_latex_results.py -- do not edit by hand unless re-export fails
\begin{table}[htbp]
\centering
\caption[Same-window mean Kendall $\tau$ by proxy family for every score.]{Same-window mean Kendall $\tau$ by proxy family for every score ($n=5{,}521$). The first three columns repeat Table~\ref{tab:alignment-hybrid} without intervals; the liquidation, inverse-risk and GMX trading-success families are archived.}
\label{tab:alignment-summary}
\fitwidth{%
\begin{tabular}{lcccccc}
\toprule
Method & Transfer & Allowance & Sybil & Liq. & Inv.risk & GMX \\
\midrule
EndorseRank & 0.330 & 0.375 & 0.231 & 0.093 & -0.055 & 0.096 \\
AWP & 0.612 & -0.024 & 0.278 & 0.141 & -0.127 & 0.236 \\
\midrule
C-PR ($\lambda=1$) & 0.333 & 0.355 & 0.233 & 0.093 & -0.056 & 0.096 \\
C-PR ($\lambda=0$) & 0.534 & -0.030 & 0.286 & 0.113 & -0.094 & 0.179 \\
C-PR ($\lambda=0.25$) & 0.531 & -0.009 & 0.288 & 0.114 & -0.093 & 0.176 \\
C-PR ($\lambda=0.5$) & 0.527 & -0.001 & 0.290 & 0.113 & -0.091 & 0.171 \\
C-PR ($\lambda=0.75$) & 0.519 & 0.006 & 0.292 & 0.110 & -0.086 & 0.164 \\
S-PR & 0.494 & 0.116 & 0.220 & 0.119 & -0.099 & 0.181 \\
\bottomrule
\end{tabular}
}
\end{table}


\dissertationSubheading[sec:parameter-summary]{매개변수 요약}

\begin{table}[htbp]
\centering
\caption{주된 평가에 쓴 기본 매개변수.}
\label{tab:parameter-summary}
\fitwidth{%
\begin{tabular}{lll}
\toprule
매개변수 & 값 & 역할 \\
\midrule
페이지랭크 감쇠 계수 $d$ & $0.85$ & 순간 이동과 섞는 기본 비율 \\
수렴 허용 오차 $\varepsilon$ & $10^{-8}$ & 거듭제곱 반복을 멈추는 기준 \\
최대 반복 수 $I_{\max}$ & $300$ & 거듭제곱 반복 횟수의 상한 \\
로지스틱 $k$ & $0.01$ & 감쇠 기울기(AWP, EndorseRank, 송금 층) \\
로지스틱 $t_0$ & $180$일 & 감쇠 중간점(AWP, EndorseRank, 송금 층) \\
값 변환 $b$ & $1$(원래 단위) & AWP와 EndorseRank의 값 무게 $V(z)$ \\
재시작 분포 & $\propto X_u$; 균등 & AWP; EndorseRank, C-PR \\
섞는 비율 $\lambda$ & $0.5$(주된 값); $0.25$, $0.75$ & C-PR의 허락 층 무게, 실행 전에 정함 \\
S-PR 재시작 벡터 & 허락 걷기, 정규화, 하한 $0$ & 보증을 씨앗으로 쓴 순간 이동 \\
홀드아웃 동결일 / 결과 기간 & 2026-02-28 / 2026-03-01 -- 2026-05-31 & 스펜더 코호트 $n=1{,}335$ \\
등록 재현 동결일 / 라벨 기간 & 2026-05-31 / 2026-06-01 -- 2026-08-31 & 스펜더·트레이더 코호트; 허용 폭 $0.02$ \\
관찰 기간 & 2025-12-01 -- 2026-05-31 & UTC, 양 끝 포함 \\
부트스트랩 다시 뽑기 수 / 시드 & $400$ / $42$ & 지갑 짝지은 부트스트랩, 백분위 95\% 신뢰구간 \\
시간 잰 반복 / 준비 실행 & $5$ / $1$ & 실행 시간 재기 절차 \\
코호트 안 부분 표본 시드 & \texttt{benchmark-scale-v1} & 늘 같게 정해지는 SHA256 시드 \\
확장 풀 부분 표본 시드 & \texttt{benchmark-tier2-v1} & 늘 같게 정해지는 SHA256 시드 \\
\bottomrule
\end{tabular}
}
\end{table}

\dissertationSubheading[sec:proxy-appendix]{대리 지표 계산 정의(참고)}

아래 공식은 \ref{ch:methodology}장의 공식이에요.

\begin{itemize}
\item 송금: $\mathrm{in\_degree}(w)=|\{u:\exists u\to w\}|$; $\mathrm{in\_value}(w)=\sum_{e:u\to w}\mathrm{value}(e)$.

\item 허락: $\mathrm{in\_approve\_degree}(w)=|\{o:\exists o\to w\}|$; $\mathrm{in\_approve\_value}(w)=\sum_{o,\mathrm{token}} a^{\star}(o,w,\mathrm{token})$. 여기서 $w$는 스펜더예요.

\item 시빌 보정 안정성: $\mathrm{inbound\_counterparty\_ratio}(w)=|\{u:u\to w\}|/|T^{\mathrm{in}}(w)|$(자기 자신에게 보낸 송금은 뺌); 날 수로 센 활동 기간; 활동한 달 수.
\end{itemize}

\dissertationSubheading[sec:glossary]{용어 풀이}

\begin{description}
\item[활발함(Activeness)] AWP에서 한 주소가 받는 재시작 무게예요. 그 주소가 송금을 보낸 주소마다 그곳으로 보낸 마지막 송금의 시간 무게를 구하고, 이것을 모두 더한 값이에요.
\item[허락(Allowance)] 스펜더가 주인 대신 토큰을 옮길 수 있도록 ERC-20에서 권한을 주는 일이에요.
\item[일치도(Alignment)] 평판 점수와 검증용 대리 지표가 순위에서 얼마나 맞는지를 말해요. 중간 순위(동점끼리 나눠 갖는 평균 순위)로 구한 스피어먼의 $\rho$와 켄달의 $\tau$로 재요. 보고한 $\tau$는 모두 동점을 보정한 $\tau_b$예요.
\item[구성개념 겹침(Construct overlap)] 순위를 매기는 방법과 대리 지표가 같은 자료에서 나오는 경우예요. 이러면 일치도가 부풀려져요.
\item[C-PR 양 끝점(C-PR end points)] $\lambda=1$일 때와 $\lambda=0$일 때의 C-PR이에요. 이 둘은 금액을 그대로 쓴 허락 층이나 송금 층 하나만을, 고르게 다시 출발하면서 걸어요. 등록 문서에서는 이 둘을 엔도스랭크와 AWP라고 불러요.
\item[보증 그래프(Endorsement graph)] 마지막 허락들로 만든 방향이 있는 그래프예요. 엔도스랭크의 입력으로 쓰이고, C-PR의 허락 층으로도 쓰여요.
\item[매칭 코호트(Matched cohort)] 주된 표본으로, 아비트럼 원에서 2025-12-01부터 2026-05-31 사이에 포지션을 적어도 세 번 닫은 GMX V2 지갑 $n=5{,}521$개예요. 허락이나 송금을 꺼내기 전에 GMX 자료로 정했어요. 어떤 지갑이 들어가는지는 그래프에서 이어져 있는지와 상관이 없어요. 이 지갑들 가운데 $966$개는 어떤 허락 화살표에도 닿지 않고, $381$개는 어떤 송금 화살표에도 닿지 않아요.
\item[확장 지갑 풀(Expanded wallet pool)] 실행 시간이 뽑은 지갑 수에 얼마나 민감한지 시험하려고 쓴 더 큰 주소 모음($N=99{,}080$)이에요. 여기에 더해 넣은 주소들은 꺼낸 화살표가 없어서, 풀이 프로그램의 그래프에 점을 하나도 더하지 않아요.
\item[짝지은 부트스트랩(Paired bootstrap)] 지갑을 복원 추출(뽑은 것을 도로 넣고 다시 뽑기)로 다시 뽑는 방법이에요. 이때 모든 계수가 번호 행렬 하나를 함께 써요. 그래서 계수 사이의 차이를 같은 다시 뽑은 표본 위에서 검사할 수 있어요.
\item[페이지랭크(PageRank)] 감쇠가 있는 무작위 걷기의 정상 분포(오래 걸었을 때 각 점에 머무는 비율)로 점에 순위를 매기는 스펙트럼 방식(행렬의 고유벡터로 순위를 매기는 방식)이에요.
\item[대리 지표(Proxy)] 순위가 그 분야의 직관과 맞는지 점검할 때 쓰는, 관찰할 수 있는 지갑 측정값이에요.
\item[스펜더 코호트(Spender cohort)] $n=1{,}335$개 지갑으로 된 홀드아웃 표본이에요. 꺼낸 허락 자료에서 $t_1$(2026년 2월 28일, 23:59:59 UTC)에 0보다 큰 마지막 허락을 적어도 하나 가지고 있던 모든 스펜더예요.
\item[순간 이동(Teleport)] 페이지랭크가 확률 $1-d$로 하는 건너뛰기예요. 재시작 분포에서 뽑은 점으로 건너뛰어요. 재시작 분포는 엔도스랭크와 C-PR에서는 고르고, AWP에서는 송금을 보낸 활동에 따라 무게를 줘요.
\end{description}

% !TEX root = ../main.tex
\section{SQL Templates, Configuration, and Extended Notes}
\label{ch:appendix-sql}

To aid reproduction, this appendix gathers representative SQL patterns, configuration keys, and extended methodological notes.

\dissertationSubheading[sec:sql-approval]{Approval log extraction pattern}

An ERC-20 \texttt{Approval} event is identified by its topic~$0$, which is the keccak256 hash of \texttt{Approval(address,address,uint256)}. For one month, a simplified query is:

\begin{verbatim}
SELECT
  block_timestamp,
  block_number,
  transaction_hash,
  log_index,
  address AS token,
  topics[SAFE_OFFSET(1)] AS owner_topic,
  topics[SAFE_OFFSET(2)] AS spender_topic,
  data AS value_data
FROM `bigquery-public-data.goog_blockchain_arbitrum_one_us.logs`
WHERE block_timestamp >= TIMESTAMP('2025-12-01')
  AND block_timestamp <  TIMESTAMP('2026-01-01')
  AND topics[SAFE_OFFSET(0)] = '0x8c5be1e5ebec7d5bd14f71427d1e84f3dd0314c0f7b2291e5b200ac8c7c3b925'
  AND (
    LOWER(CONCAT('0x', SUBSTR(topics[SAFE_OFFSET(1)], 27))) IN UNNEST(@wallet_list)
    OR LOWER(CONCAT('0x', SUBSTR(topics[SAFE_OFFSET(2)], 27))) IN UNNEST(@wallet_list)
  );
\end{verbatim}

With rows filtered against \texttt{@wallet\_list}, scan volume is proportional to the cohort instead of the full chain. Monthly batches keep every query below the per-query byte caps documented in the study's acquisition record.

\dissertationSubheading[sec:sql-transfer]{Transfer log extraction pattern}

\texttt{Transfer} events are selected by the topic~$0$ of \texttt{Transfer(address,address,uint256)}:

\begin{verbatim}
SELECT
  block_timestamp,
  block_number,
  transaction_hash,
  log_index,
  address AS token,
  topics[SAFE_OFFSET(1)] AS from_topic,
  topics[SAFE_OFFSET(2)] AS to_topic,
  data AS value_data
FROM `bigquery-public-data.goog_blockchain_arbitrum_one_us.logs`
WHERE block_timestamp >= TIMESTAMP('2025-12-01')
  AND block_timestamp <  TIMESTAMP('2026-01-01')
  AND topics[SAFE_OFFSET(0)] = '0xddf252ad1be2c89b69c2b068fc378daa952ba7f163c4a11628f55a4df523b3ef'
  AND (
    LOWER(CONCAT('0x', SUBSTR(topics[SAFE_OFFSET(1)], 27))) IN UNNEST(@wallet_list)
    OR LOWER(CONCAT('0x', SUBSTR(topics[SAFE_OFFSET(2)], 27))) IN UNNEST(@wallet_list)
  );
\end{verbatim}

Before PageRank, post-processing keeps, for each (token, owner, spender) triple, the value of the last \texttt{Approval} event in (\texttt{block\_number}, \texttt{log\_index}) order and drops zero (revoked) values. It then weights each latest approval and each transfer by its time weight and bounded value weight, builds the raw-amount layers of C-PR beside them, and keeps every edge with at least one endpoint in the matched cohort. Amounts stay in raw token base units, the \texttt{uint256} value cast to a float, and no decimal scaling is applied before the value weight or, in the C-PR layers, before allowances or transfer values are summed across tokens.

\dissertationSubheading[sec:config-keys]{Configuration keys}

All of the pipeline's parameters are kept in a single versioned configuration file, whose principal keys include:

\begin{itemize}
\item start and end timestamps of the observation window;
\item the PageRank damping factor, tolerance, and maximum number of iterations;
\item logistic decay parameters $k$ and $t_0$ and the value-weight parameter $b$;
\item subsample seeds and the number of benchmark repetitions;
\item the locations to which intermediate artifacts and table fragments are written (\texttt{results/tables/}).
\end{itemize}

Changing the damping factor or the decay parameter requires a rerun of the entire evaluation, including the robustness sweeps and the table export. All figures must then be regenerated from the same artifacts so that the reported numbers match the pipeline's machine-readable summary.

\dissertationSubheading[sec:complexity-notes]{Complexity notes}

Write $m=|E|$. Building the latest-allowance edges sorts the $R_a$ approval rows in (\texttt{block\_number}, \texttt{log\_index}) order, in $O(R_a\log R_a)$ time, and the transfer edges take $O(R_t)$ over the transfer rows. PageRank over $I$ iterations costs $O(m+mI)$, the first term for building the sparse operator, which dominates the measured calls (Appendix~\ref{sec:pr-sparsity}). The edge lists and score vectors account for most of the memory, so the peak-memory ratio in Table~\ref{tab:benchmark-runtime} is of the order of the edge ratio.

\dissertationSubheading[sec:interpretation-checklist]{Conventions for the result tables}

The tables of Chapter~\ref{ch:results} follow these conventions:

\begin{enumerate}
\item Every $\tau$ is Kendall's $\tau_b$ on raw scores and proxies, reported with its 95\% paired-bootstrap interval. Differences between coefficients come from the paired contrast tables, not from subtracting point estimates.
\item On the matched cohort, EndorseRank's coefficients are capped by its ties (Section~\ref{sec:rank-divergence}), so they are compared with other coefficients only through the contrasts and with that cap in mind.
\item A runtime is the time of one PageRank call on edge lists built in advance, covering matrix assembly and power iteration but not extraction or preprocessing. The same graph under the same settings shows the same numbers in every table (Section~\ref{sec:computational-benchmark}).
\item The sample-definition sweep (Table~\ref{tab:robustness-sample-size}) is a sensitivity analysis, and its absolute values are not compared with the matched-cohort tables.
\end{enumerate}

\dissertationSubheading[sec:archived-baselines]{Archived extended baselines}

The archived matrices, and the reasons they stay out of the main comparison, are described in Appendix~\ref{ch:appendix-archive}.

\dissertationSubheading[sec:limitations-reprise]{Limitations}

The SQL above was run only on Arbitrum One. All reported numbers come from the pipeline's machine-readable summaries for the matched cohort ($n=5{,}521$), the pool subsamples, the spender cohort ($n=1{,}335$) and the exploratory trader run; the registered June--August results come from the summary of that run (Section~\ref{sub:fresh-results}).

% !TEX root = ../main.tex
\section{페이지랭크와 순위 상관의 기본 성질}
\label{ch:appendix-theory}

\ref{ch:literature}--\ref{ch:results}장은 여기에 정리한 기본 성질에 기대고 있어요. 이 내용은 널리 알려진 표준적인 것이에요. 페이지랭크(PageRank)의 성질은 Page 등\cend{15}과 Langville과 Meyer\cend{23}를 따르고, 순위 상관(두 줄 세우기가 얼마나 비슷한지 재는 수)은 Spearman\cend{20}과 Kendall\cend{21}을 따라요.

\dissertationSubheading[sec:pr-existence]{페이지랭크의 존재와 유일성}

$P$를 점 $|V|$개 위의 행 확률 행렬(각 행의 합이 1인 행렬)이라고 해요. 필요하면 식~\eqref{eq:transition}의 방법으로 매달린 점(나가는 화살표가 없는 점)을 고친 다음의 행렬이에요. $d\in(0,1)$일 때 식~\eqref{eq:google-matrix}의 구글 행렬은 다음과 같아요.
\[
G_{\mathrm{PR}}=d P^{\top}+\frac{1-d}{|V|}\mathbf{1}\mathbf{1}^{\top},
\]
이 행렬은 모든 칸이 양수이고, 각 열의 합이 1인 열 확률 행렬이에요. 양수 행렬에 대한 페론--프로베니우스(Perron--Frobenius) 정리에 따르면, $G_{\mathrm{PR}}$에는 크기가 1인 고윳값(행렬이 벡터를 몇 배로 늘리는지 나타내는 수)이 하나뿐이에요. 그리고 그에 맞는 고유벡터 $\mathbf{r}$는 모든 칸이 0보다 커요. 우리는 이 벡터를 $\sum_i r_i=1$이 되도록 크기를 맞출 수 있어요. 따라서 감쇠 계수(화살표를 따라갈 확률)와 매달린 점 규칙을 정하고 나면, 방향이 있는 그래프는 모두 단 하나의 페이지랭크 점수를 가져요\cend{23}.

이 논문의 재시작 벡터 가운데 두 개는 고르지 않아요. AWP(적응형 가중 페이지랭크)의 재시작 벡터는 활발함(활동 정도)에 비례하므로, 아무것도 보내지 않은 주소에서는 0이에요. S-PR(보증 씨앗 페이지랭크)의 재시작 벡터는 허락 걷기의 점수를 합이 1이 되게 맞춘 것이므로, 허락 점수가 없는 송금 층 지갑에서는 0이에요. 그래서 두 구글 행렬 모두 모든 칸이 양수이지는 않아요. 그래도 고정점(계산을 되풀이해도 더는 바뀌지 않는 값)은 여전히 하나뿐이에요. 두 확률 벡터의 차이를 보면 재시작 항이 서로 지워지고, 연산자는 두 벡터 사이의 $\|\cdot\|_1$ 거리를 많아야 $d$배로 만들기 때문이에요. 하지만 일반적으로 이 고정점의 모든 칸이 0보다 크지는 않아요. 재시작 몫이 없고, 재시작 지갑에서 출발한 어떤 길도 닿지 않는 지갑은 점수가 계속 0으로 남기 때문이에요.

답이 하나뿐이라는 성질(유일성) 덕분에 엔도스랭크(EndorseRank)와 AWP를 분명하게 비교할 수 있어요. 매개변수를 정하고 나면, 두 방법은 저마다 주어진 화살표 묶음에 꼭 하나의 점수 벡터를 줘요. 그래서 $\mathbf{s}_{\mathrm{ER}}$와 $\mathbf{s}_{\mathrm{AWP}}$의 차이는 모두 두 방법의 화살표와 재시작 분포에서 생길 수밖에 없어요. 여러 풀이 프로그램 탓으로 돌릴 수는 없어요.

\dissertationSubheading[sec:pr-continuity]{화살표 무게에 대한 연속성}

$d\in(0,1)$를 고정하면, 페이지랭크는 $P$의 칸 값을 따라 연속적으로(끊김 없이) 바뀌어요. $\tilde P$가 같은 점들 위의 또 다른 행 확률 행렬이고 그 페이지랭크 벡터가 $\tilde{\mathbf{r}}$라면, 다음 부등식이 맞아요.
\[
\|\mathbf{r}-\tilde{\mathbf{r}}\|_1\le\frac{d}{1-d}\,\max_i\bigl\|P_{i\cdot}-\tilde P_{i\cdot}\bigr\|_1,
\]
이 한계는 $|V|$에 달려 있지 않아요\cend{23}.

$P$에 대한 연속성은 화살표 무게에 대한 연속성과 같지 않아요. 행 $P_{i\cdot}=w_{i\cdot}/o_i$는 나가는 무게의 합 $o_i$가 0보다 큰 동안에만 무게를 따라 연속적으로 움직여요. 어떤 점이 처음으로 나가는 화살표를 얻거나 마지막 나가는 화살표를 잃으면, 그 행은 고른 매달린 행과 나가는 이웃에게 몰린 행 사이를 훌쩍 건너뛰어요. 관련된 무게가 아무리 작아도 그래요. \ref{sec:pagerank-worked-example}절의 계산 예에서, 무게가 $10^{-12}$인 화살표 $D\to A$를 더하면 $\mathbf{r}$가 $\|\cdot\|_1$ 노름(칸마다 차이의 크기를 모두 더한 값)으로 $0.37$만큼 움직이고, $r_A$는 $0.1375$에서 $0.3202$로 올라가요.

감쇠 계수를 바꿔 보는 실행은 연속성으로 설명돼요. 이 실행은 $P$를 그대로 두고 $d$만 매끄럽게 바꾸는데, $d<1$이면 페이지랭크 벡터는 $d$를 따라 연속적으로 바뀌기 때문이에요. 상위 토큰(가장 많이 쓰인 토큰)만 남기는 거르기는 연속성으로 설명되지 않아요. 이 거르기는 화살표를 지우는데, 그러면 어떤 지갑은 나가는 화살표가 하나도 없게 될 수 있어요. 그래서 이 점검이 얼마나 안정적인지는 오로지 자료에 달려 있어요(표~\ref{tab:robustness-tokens}). 실제 자료에서 엔도스랭크의 허락 지표 가족(비슷한 지표 묶음) 평균은 $d\in\{0.75,0.85,0.95\}$에서 소수점 아래 세 자리까지 바뀌지 않았어요(표~\ref{tab:robustness-damping}). 코호트의 대부분이 동점이니, 순위는 거의 움직일 수 없었을 거예요.

\dissertationSubheading[sec:pr-sparsity]{희소성과 반복 한 번의 비용}

0이 아닌 화살표가 $m=|E|$개 있으면, 희소한(듬성듬성한) 곱 $P^{\top}\mathbf{x}$에는 $O(m)$번의 셈이 들어요. 화살표 목록으로 희소 연산자를 만드는 것도 $O(m)$이에요. 그래서 $I$번 하는 거듭제곱 반복(같은 계산을 되풀이하기)은 $O(m+mI)$예요. 매칭 코호트(모든 점수가 똑같이 점수를 매기는 지갑 묶음)에서 송금 그래프는 보증 그래프보다 화살표가 약 $9.3$배 많아요. 그래서 $I$가 똑같더라도 송금 방법은 대략 한 자릿수(약 열 배)만큼 더 오래 걸릴 거예요. 실제로 잰 호출에서는 연산자를 조립하는 데 시간이 대부분 들고, 반복에는 조금만 들어요(\ref{sec:benchmark-results}절). 그래서 잰 시간의 비율 $9.2$는 화살표 수의 비율을 따라가요.

\dissertationSubheading[sec:spearman-properties]{스피어먼의 \texorpdfstring{$\rho$}{rho}의 성질}

스피어먼의 로($\rho$)는 순위로 계산한 피어슨 상관(두 값이 함께 커지고 작아지는 정도)이에요. 이 값은 $[-1,1]$ 안에 있어요. 그리고 두 변수 가운데 어느 쪽을 늘 커지기만 하는 방식으로 바꾸어도 달라지지 않아요. $\rho=1$이면 순서가 완전히 같은 방향으로 맞고, $\rho=-1$이면 완전히 거꾸로 맞는다는 뜻이에요. 같은 값(동점)에는 중간 순위(동점끼리 나눠 갖는 평균 순위)를 줘요. 평가 처리 과정은 동점을 제대로 다루는 표준 구현을 써요. 평판 점수와 대리 지표(대신 재는 값)는 연속적인 값이고 꼬리가 두꺼워서(아주 큰 값이 드물게 나와서), 원래 값에 대한 피어슨 상관보다 순위에 바탕을 둔 관계 재기가 더 잘 맞아요. AWP 기준선(비교 기준) 논문도 마찬가지로 자기 순위를 순위 통계로 평가해요\cend{3}.

\dissertationSubheading[sec:kendall-properties]{켄달의 \texorpdfstring{$\tau$}{tau}의 성질}

켄달의 타우($\tau$)는 순서가 맞는 쌍의 수 $n_c$와 순서가 어긋나는 쌍의 수 $n_d$를 견주어요. 동점 보정을 하지 않으면, $n_0=n(n-1)/2$일 때 $\tau_a=(n_c-n_d)/n_0$는 무작위로 뽑은 두 지갑을 두 순위가 같은 방향으로 줄 세울 확률에서 반대 방향으로 줄 세울 확률을 뺀 값이에요\cend{21}. 보고한 $\tau$는 모두 식~\eqref{eq:kendall}의 동점 보정 $\tau_b$예요. 동점이 없으면 두 값은 같지만, 여기서는 동점이 아주 많아요. 엔도스랭크는 매칭 코호트에서 서로 다른 값을 $99$개만 가져요. 그래서 모든 지갑 쌍의 $67.3\%$가 엔도스랭크에서 동점인데, AWP에서는 $0.5\%$뿐이에요. 그리고 엔도스랭크와 AWP 사이에 보고한 $\tau_b=0.360$은 $\tau_a=0.206$과 짝이 돼요. 동점은 $\tau_b$의 위쪽 한계도 정해요. 엔도스랭크가 얻을 수 있는 가장 큰 값은 허락 대리 지표에 대해서는 약 $0.38$이에요. 허락 대리 지표는 엔도스랭크가 떼어 내는 바로 그 131개 지갑에서만 0보다 커요. 동점이 없는 대리 지표에 대해서는 약 $0.57$이에요. 그래서 허락 가족에서의 $0.375$(표~\ref{tab:alignment-hybrid})는 두 순위가 떼어 내는 쌍에서는 거의 완전히 맞는다는 뜻이에요. 그리고 이 값은 천장(얻을 수 있는 가장 큰 값)이 다른 계수와 곧바로 비교할 수 없어요.

$n$이 크면 $\tau$와 $\rho$는 보통 부호가 같지만, 크기는 달라요. 엔도스랭크나 AWP를 대리 지표와 짝지은 것 가운데 $|\tau|>0.02$인 짝들에서, 비율 $\rho/\tau$는 $1.01$과 $1.49$ 사이에 있고 늘 $\rho$가 더 커요. AWP 기준선이 그렇게 하듯이, 우리도 둘 다 보고해요\cend{3}. 우리의 주된 요약 통계는 짝지은 부트스트랩(다시 뽑기) 구간을 붙인 가족 평균 $\tau$예요\cend{34}.

\dissertationSubheading[sec:construct-validity-note]{구성 타당도에 대한 메모}

Cronbach와 Meehl\cend{22}은 기준 관련 타당도(정답 기준과 얼마나 맞는지)와 구성 타당도(점수가 재려는 것을 정말 재는지)를 나누어요. 확실한 정답 기준이 되는 채무 불이행(빌린 돈을 갚지 못함) 라벨이 없으므로, 이 연구는 구성개념(점수가 재려고 하는 생각)을 중심에 둔 방법을 써요. 각 점수는 그 점수의 개념 정의를 담은 검증 점검으로 시험해요(예를 들어 엔도스랭크는 허락을 받는 것, AWP는 들어오는 송금을 받는 것). 그리고 두 라벨(나중에 채점에 쓰는 정답 값), 곧 새 허락 쌍과 새 송금자가 있는 홀드아웃(떼어 두었다가 나중에 맞혀 보는) 기간으로도 시험해요. 같은 구성개념 안에서 $\tau$가 높으면 이 해석을 뒷받침해요. 반대로 서로 다른 구성개념 사이에서 $\tau$가 낮다고 해서 저절로 실패를 뜻하지는 않아요. 그것은 판별 타당도(서로 다른 것을 다르게 재는지)를 보여 주는 것일 수도 있어요.

\dissertationSubheading[sec:numerical-summary-sheet]{주요 값을 보고한 곳}

아래 목록은 주요 값마다 그 값을 내놓은 표를 알려 줘요.

\begin{center}
\fitwidth{%
\begin{tabular}{ll}
\toprule
값 & 보고한 곳 \\
\midrule
매칭 코호트 크기, $|V|$, $|E|$, 실행 시간, 표준편차, 메모리, 반복 수(방법 네 개) & 표~\ref{tab:benchmark-runtime} \\
코호트 안 규모 키우기 단계 & 표~\ref{tab:benchmark-scaling-incohort} \\
풀 부분 표본 단계($N=99{,}080$) & 표~\ref{tab:benchmark-scaling} \\
95\% 신뢰구간을 붙인 가족 평균 $\tau$, 모든 점수 & 표~\ref{tab:alignment-hybrid} \\
C-PR과 S-PR의 짝지은 $\Delta\tau$ 대비 & 표~\ref{tab:tau-diff-hybrid} \\
방법 사이 $\tau$; 그래프 밖 지갑을 외톨이 점으로 둔 경우 & 표~\ref{tab:tie-sensitivity} \\
감쇠 계수, 상위 토큰, 표본 정의 바꿔 보기 & 표~\ref{tab:robustness-damping}, \ref{tab:robustness-tokens}, \ref{tab:robustness-sample-size} \\
스펜더 홀드아웃($n=1{,}335$) & 표~\ref{tab:holdout-spenders}, \ref{tab:holdout-diff}와~\ref{tab:holdout-receiving} \\
탐색적 트레이더 실행 & 표~\ref{tab:holdout-traders} \\
등록한 6월부터 8월까지의 재현 & \ref{sub:fresh-results}절 \\
EndorseRank와 AWP 비교, 대리 지표별 결과 & 부록~\ref{ch:appendix-er-awp} \\
관찰 기간 & 2025-12-01 -- 2026-05-31(표~\ref{tab:parameter-summary}) \\
\bottomrule
\end{tabular}
}
\end{center}

여기 적은 값은 모두 평가 요약에도 있고, \ref{ch:results}장이나 부록의 표에도 나와요.

% !TEX root = ../main.tex
\section{Archived Extended Baselines}
\label{ch:appendix-archive}

Further PageRank variants were scored on the same matched cohort and proxy families. They are not among the scores evaluated in Chapter~\ref{ch:results}.

\dissertationSubheading[sec:seven-method-archive]{Seven-method alignment matrix}

The evaluation archive keeps the seven-method matrices (LP-PR, CW-AWP, LF-PR, and RiskProp-style variants), which are not drawn here. Their rows were computed with the two raw-amount layers of C-PR walked alone as the allowance and transfer references. Among these methods, the allowance walk leads on the allowance family, and the transfer walk, together with its closely related variants, leads on transfer.

The archive matrices show the following patterns:

\begin{itemize}
\item Methods that use transfer-like edges cluster around the transfer walk on the transfer and Sybil families.
\item When a method has very few edges, its rankings can be unstable or degenerate, and a low runtime may simply reflect a tiny $m$.
\item RiskProp-style rows that match the transfer walk numerically indicate shared transfer structure; they do not constitute an independent discovery.
\end{itemize}

\dissertationSubheading[sec:six-aave-archive]{Six-method Aave-oriented matrix}

An Aave-oriented preset is archived as well. Within the study window, several of the lending-event graphs on Arbitrum are sparse, and the delegation-based rows may be degenerate when \texttt{BorrowAllowanceDelegated} events were not extracted (Appendix~\ref{ch:appendix-eval}).

That preset remains in the archive and is not drawn here.

Although these matrices motivate future work on credit-delegation edges and cross-protocol lending labels, they leave Claims C1--C5 unchanged.

\dissertationSubheading[sec:why-archived]{Why these variants are archived}

Three reasons explain why the scores evaluated are EndorseRank and AWP, with C-PR and S-PR built from the same two layers:

\begin{enumerate}
\item Allowance-based endorsement is the new construct, and evaluating it beside the established transfer method~\cend{3} isolates the research question.
\item Variants that share transfer edges with AWP contribute extra rows but no independent construct.
\item The outcome-native GF-PR reference from the research proposal is derived from the same GMX events as the trading-success proxies, so any alignment between the two counts as construct overlap and not as validation (Section~\ref{sub:gf-pr}). The GF-PR three-method matrix is archived alongside the seven-method and six-Aave presets.
\end{enumerate}

\dissertationSubheading[sec:archive-repro]{Reproducing archive tables}

Running the evaluation pipeline's table export with the seven-method and six-Aave presets regenerates the archive tables. CSV copies are written out together with the table fragments.

\dissertationSubheading[sec:archive-method-notes]{Notes on archived method identifiers}

The archive tables and CSV exports use the identifiers below.

\begin{description}
\item[LP-PR] The raw-amount transfer layer, with AWP's time decay, with each edge weight multiplied by one minus the GMX liquidation rate of its sender and one minus that of its receiver. It tracks the transfer walk closely on the transfer families. Its weights use the same liquidation outcomes as the liquidation proxies, so it shares the circularity of GF-PR.
\item[CW-AWP] The raw-amount transfer layer restricted to transfers of a fixed list of collateral tokens. Its transfer $\tau$ is slightly lower than that of the full layer.
\item[LF-PR] A lending-flow graph from Aave events, with borrows as edges from the pool to the user and repayments and liquidations as edges from the user to the pool, all under AWP's time decay. Few cohort wallets used Aave on Arbitrum in the window, so the graph has very few edges and near-zero runtimes, and its rankings cannot be compared with the transfer walks.
\item[RiskProp] The raw-amount transfer layer with each edge weight multiplied by one minus the GMX liquidation rate of its sender. It borrows the idea of propagating risk from Lin et al.~\cend{12} but is not their method. On this cohort it mirrors the transfer walk numerically on several families, so the liquidation weighting changed little.
\item[LiqCall-PR / Borrow-PR / Repay-PR] Aave-oriented wallet-to-wallet graphs derived from liquidation calls, borrows, and repays, where applicable. On Arbitrum in the study window they have only tens to hundreds of edges, which leaves alignment weak and runtimes negligible.
\item[Delegation-PR] Intended to use credit-delegation allowances. With zero delegation edges extracted, its scores are degenerate and its $\tau$ cannot be interpreted.
\end{description}

No archive row is a peer-reviewed baseline.

\dissertationSubheading[sec:archive-decision-log]{Criteria for the scores evaluated}

Chapters~\ref{ch:results}--\ref{ch:discussion} concentrate on EndorseRank, AWP and the walks that couple their two layers; the other variants fall short of the following criteria:

\begin{enumerate}
\item Construct independence: a baseline should implement different edge semantics instead of merely reweighting transfers.
\item Data adequacy: there should be enough edges for PageRank to be stable on the matched cohort.
\item Scope: the contribution is allowance-based endorsement; it is not meant to survey every possible PageRank variant.
\end{enumerate}

By these criteria, EndorseRank (allowance) and AWP (transfer), with the walks that couple their layers, make up a minimal sufficient set. The seven-method and six-Aave matrices remain archived and do not enter into Claims~C1--C5.


% !TEX root = ../main.tex
\section{엔도스랭크와 AWP 나란히 보기}
\label{ch:appendix-er-awp}

\ref{ch:results}장은 이 논문의 모든 점수를 함께 보고하고, 각 페이지랭크(PageRank)를 그것이 매끄럽게 다듬는(평활하는) 원시 차수(화살표를 그냥 센 수)와 견주어 읽어요. 이 부록은 한 층만 걷는 두 점수, 곧 엔도스랭크(EndorseRank)와 AWP(적응형 가중 페이지랭크) 사이의 직접 비교를 모아 두었어요. 같은 기간에서의 가족 평균(비슷한 지표 묶음에 대한 평균값)과 짝지은 차이, 대리 지표(대신 재는 값)별 계수, 그리고 기간 밖에서의 짝지은 차이예요. 이 비교들은 하나도 미리 정해 두지 않았고, 어느 것도 이 논문의 주장을 담지 않아요. 두 점수는 서로 다른 화살표 위를 걷고, 서로 다른 구성개념(점수가 재려고 하는 생각)을 위해 만들어졌어요. 그래서 점검마다 앞서는 점수는 거의 언제나, 그 점검이 세는 화살표를 쓰는 점수예요. 두 점수가 다루는 범위도 달라요. 매칭 지갑 가운데 966개에는 허락 화살표가 하나도 닿지 않고, 381개에는 송금 화살표가 하나도 닿지 않아요. 또 홀드아웃(나중 기간을 떼어 두고 맞혀 보는 시험) 스펜더 1,335개 가운데 318개는 동결일(점수를 고정한 날) 전에 송금 화살표가 없어요. 같은 기간에서 엔도스랭크의 순서는 대부분 동점이 정해요(\ref{sec:rank-divergence}절). 두 홀드아웃 기간 모두에서 두 점수는 저마다 자기가 다듬는 원시 차수보다 낮아요(\ref{sec:holdout-results}절). 그래서 둘 사이에 차이가 있다고 해서, 어느 쪽이든 자기 차수에 무언가를 더한다는 것을 보여 주지는 않아요.

\dissertationSubheading[sec:er-awp-same-window]{같은 기간}

표~\ref{tab:alignment-hybrid}는 두 점수의 가족 평균을 구간과 함께 보여 주고, 표~\ref{tab:tie-sensitivity}의 마지막 행은 두 점수가 서로 얼마나 맞는지 보여 줘요. 두 모습은 엇갈려요. AWP는 송금 가족을 가장 가깝게 따라가고($0.612$) 허락 가족에서는 조금 음수예요. 반면 엔도스랭크는 허락 가족을 가장 가깝게 따라가고($0.375$), 송금 가족은 중간 정도로 따라가요($0.330$). 표~\ref{tab:tau-diff}는 같은 지갑 다시 뽑기 표본 위에서 짝지은 차이 다섯 개를 검사해요(\ref{sec:statistical-tests}절). 이 차이들에는 표의 행 순서대로 i부터 v까지 번호를 붙였어요.

% Real BigQuery data -- regenerate with:
%   2-Dissertation-Draft-Works/wallet-reputation-experiments/scripts/run_dissertation_eval.py --real --export-latex
% Generated by export_latex_results.py -- do not edit by hand unless re-export fails
\begin{table}[htbp]
\centering
\caption[Paired bootstrap tests of differences between Kendall $\tau$ coefficients.]{Paired bootstrap tests of differences between Kendall $\tau$ coefficients (matched cohort, $n=5{,}521$). Family entries use the family-mean $\tau$; the same wallet resamples are used for $a$ and $b$, so the interval is on the paired difference. An interval excluding 0 indicates a difference not attributable to sampling variation at the 5\% level. CI = 95\% percentile bootstrap over wallets (400 paired resamples, seed 42).}
\label{tab:tau-diff}
\fitwidth{%
\begin{tabular}{lccccc}
\toprule
Contrast ($a$ minus $b$) & $\tau_a$ & $\tau_b$ & $\Delta\tau$ & 95\% CI of $\Delta\tau$ & $P(\Delta\tau>0)$ \\
\midrule
EndorseRank allowance minus EndorseRank transfer & 0.375 & 0.330 & +0.045 & [-0.001, 0.084] & 0.970 \\
AWP transfer minus EndorseRank transfer & 0.612 & 0.330 & +0.282 & [0.263, 0.299] & 1.000 \\
EndorseRank allowance minus AWP allowance & 0.375 & -0.024 & +0.398 & [0.356, 0.432] & 1.000 \\
AWP sybil-stability minus EndorseRank sybil-stability & 0.278 & 0.231 & +0.047 & [0.031, 0.062] & 1.000 \\
EndorseRank in-approve degree minus EndorseRank in-degree & 0.376 & 0.341 & +0.035 & [-0.015, 0.079] & 0.925 \\
\bottomrule
\end{tabular}
}
\end{table}


다섯 대비(짝지은 차이) 가운데 셋에서 구간이 0을 벗어나요. AWP의 송금 일치도가 엔도스랭크보다 높고($+0.282$), 엔도스랭크의 허락 일치도가 AWP보다 높아요($+0.398$). 그리고 시빌(가짜 주소를 많이 만들어 속이는 공격) 안정성 가족에서는 AWP가 엔도스랭크보다 앞서요($+0.047$). 나머지 둘은 엔도스랭크를 자기 자신과 비교하는데, 둘 다 차이가 보이지 않아요. 엔도스랭크의 허락 일치도가 송금 일치도보다 높다는 것은 보이지 않았어요(대비~i, $+0.045$, 구간 $[-0.001,0.084]$). 받은 허락 수 일치도가 입차수 일치도보다 높다는 것도 보이지 않았어요(대비~v, $+0.035$, $[-0.015,0.079]$). 이 두 ‘차이 없음’은 어느 쪽의 증거로도 읽을 수 없어요. 엔도스랭크의 동점 때문에 $\tau_b$의 천장(얻을 수 있는 가장 큰 값)은 허락 대리 지표에서는 약 $0.38$이고, 동점이 없는 대리 지표에서는 약 $0.57$이에요(\ref{sec:same-window-alignment}절). 그래서 이 계수들은 같은 잣대 위에 있지 않아요. 허락 그래프 밖의 지갑에 0점을 주는 대신 외톨이 점(화살표가 하나도 없는 점)으로 남겨 두면, 대비~i와~v는 $+1.061$과 $+1.090$이 되고 구간도 0보다 훨씬 위에 있어요.

대비~ii와~iii의 부호가 반대인 것은 구성개념에서 나와요. 각 점수는 자기 화살표로 만든 가족에서 앞서요. 그리고 각 가족의 주된 대리 지표에는 그에 맞는 페이지랭크가 다듬는 입차수가 들어 있어요. 허락 가족에서 엔도스랭크가 앞서는 것은 허락을 받는 지갑 131개 덕분이에요. 엔도스랭크는 이 지갑들을 모두 공통 값보다 위로 올리지만, AWP는 이 131개를 자기 순위의 대부분에 걸쳐 흩어 놓아요. 송금 가족에서 AWP가 앞서는 것은 엔도스랭크의 송금 계수와 견준 결과인데, 이 계수는 빠진 지갑을 외톨이 점으로 남겨 두면 음수로 바뀌어요. 두 차이 모두 어느 점수가 평판을 더 잘 재는지는 말해 주지 않아요.

표~\ref{tab:robustness-sample-size}의 단계별 풀 부분 표본은 그 값을 매칭 코호트의 값과 비교할 수 없어요. 이 부분 표본들에서 각 가족 안의 두 점수 순서는 모든 단계에서 같아요. 송금 가족에서는 AWP가 앞서고, 허락 가족에서는 엔도스랭크가 앞서요.

\dissertationSubheading[sec:proxy-family-detail]{대리 지표별 자세한 결과}

가족 평균 하나는 대리 지표별 계수 여러 개를 하나로 줄인 거예요. 표~\ref{tab:alignment-transfer}부터~\ref{tab:alignment-sybil-stability}까지는 엔도스랭크와 AWP에 대해 이 계수들을 대리 지표 하나씩 보여 줘요. 스피어먼의 로($\rho$), 켄달의 타우($\tau$, 두 줄 세우기가 얼마나 같은 순서인지 나타내는 수), 그리고 $\tau$의 95\% 구간이에요. 상관을 구하기 전에, 모든 대리 지표는 값이 클수록 기대되는 평판이 높다는 뜻이 되도록 방향을 맞춰요. 송금 대리 지표는 AWP를 처음 낸 논문이 AWP를 검증할 때 쓴 바로 그 지표들이에요\cend{3}.

\dissertationSubsubheading[sub:transfer-family]{송금 가족}

% Real BigQuery data -- regenerate with:
%   2-Dissertation-Draft-Works/wallet-reputation-experiments/scripts/run_dissertation_eval.py --real --export-latex
% Generated by export_latex_results.py -- do not edit by hand unless re-export fails
\begin{table}[htbp]
\centering
\caption[Domain alignment with transfer proxies.]{Domain alignment with transfer proxies (inbound in-degree and in-value, the validation proxies used for the AWP baseline). Kendall $\tau$ and Spearman $\rho$ between reputation scores and proxy rankings. CI = 95\% percentile bootstrap over wallets (400 paired resamples, seed 42).}
\label{tab:alignment-transfer}
\fitwidth{%
\begin{tabular}{lccc}
\toprule
Proxy & Spearman $\rho$ & Kendall $\tau$ & 95\% CI of $\tau$ \\
\midrule
\multicolumn{4}{l}{\textit{EndorseRank}} \\
In-degree & 0.418 & 0.341 & [0.319, 0.363] \\
In-value (sum inbound) & 0.396 & 0.319 & [0.300, 0.340] \\
\addlinespace
\multicolumn{4}{l}{\textit{AWP}} \\
In-degree & 0.884 & 0.728 & [0.718, 0.737] \\
In-value (sum inbound) & 0.669 & 0.495 & [0.481, 0.510] \\

\bottomrule
\end{tabular}
}
\end{table}


표~\ref{tab:alignment-transfer}에서, AWP와 입차수(들어오는 화살표 수)의 일치는 매칭 코호트에서 대리 지표별 켄달의 $\tau$ 가운데 가장 커요. 받은 양의 합과의 일치는 그보다 훨씬 낮은데, 그 까닭은 \ref{sec:same-window-alignment}절에 있어요. 엔도스랭크의 송금 계수는 중간 정도이고($0.341$과 $0.319$), 그 구간은 0과도, AWP의 구간과도 겹치지 않아요. 이 계수들은 주로, 허락 그래프 안의 지갑이 그 밖의 지갑보다 송금도 더 많이 받는다는 점을 보여 줘요(\ref{sec:rank-divergence}절). 평판을 들어오는 흐름과 같은 것으로 보는 서비스라면, 그 흐름은 입차수만 봐도 읽을 수 있어요. 그리고 기간 밖에서 AWP는 새 송금자를 입차수보다 못하게 줄 세워요(\ref{sec:holdout-results}절).

\dissertationSubsubheading[sub:allowance-family]{허락 가족}

% Real BigQuery data -- regenerate with:
%   2-Dissertation-Draft-Works/wallet-reputation-experiments/scripts/run_dissertation_eval.py --real --export-latex
% Generated by export_latex_results.py -- do not edit by hand unless re-export fails
\begin{table}[htbp]
\centering
\caption[Domain alignment with allowance proxies.]{Domain alignment with allowance proxies (spender-side in-approve degree and value). CI = 95\% percentile bootstrap over wallets (400 paired resamples, seed 42).}
\label{tab:alignment-allowance}
\fitwidth{%
\begin{tabular}{lccc}
\toprule
Proxy & Spearman $\rho$ & Kendall $\tau$ & 95\% CI of $\tau$ \\
\midrule
\multicolumn{4}{l}{\textit{EndorseRank}} \\
In-approve degree & 0.380 & 0.376 & [0.344, 0.403] \\
In-approve value (sum) & 0.380 & 0.373 & [0.341, 0.398] \\
\addlinespace
\multicolumn{4}{l}{\textit{AWP}} \\
In-approve degree & -0.029 & -0.024 & [-0.043, -0.002] \\
In-approve value (sum) & -0.029 & -0.023 & [-0.043, -0.002] \\

\bottomrule
\end{tabular}
}
\end{table}


표~\ref{tab:alignment-allowance}에서 엔도스랭크의 두 계수는 거의 같고, 동점 때문에 얻을 수 있는 가장 큰 값(천장)에 닿아 있어요(\ref{sec:same-window-alignment}절). AWP의 두 계수는 조금 음수이고($-0.024$와 $-0.023$), 구간은 0을 벗어나요. 곧 코호트 지갑들 가운데서, 송금 그래프에서 중심에 있는 정도(중심성)로는 스펜더로 허락받은 지갑을 골라내지 못해요. 이 비교는 허락을 받는 지갑 131개에 기대고 있고, 처음부터 어느 정도 생길 수밖에 없는 결과예요. 엔도스랭크가 바로 그 허락들로 계산되기 때문이에요.

\dissertationSubsubheading[sub:sybil-family]{시빌 보정 안정성 가족}

% Real BigQuery data -- regenerate with:
%   2-Dissertation-Draft-Works/wallet-reputation-experiments/scripts/run_dissertation_eval.py --real --export-latex
% Generated by export_latex_results.py -- do not edit by hand unless re-export fails
\begin{table}[htbp]
\centering
\caption[Domain alignment with Sybil-adjusted stability proxies from transfer activity.]{Domain alignment with Sybil-adjusted stability proxies from transfer activity. CI = 95\% percentile bootstrap over wallets (400 paired resamples, seed 42).}
\label{tab:alignment-sybil-stability}
\fitwidth{%
\begin{tabular}{lccc}
\toprule
Proxy & Spearman $\rho$ & Kendall $\tau$ & 95\% CI of $\tau$ \\
\midrule
\multicolumn{4}{l}{\textit{EndorseRank}} \\
Inbound counterparty ratio & 0.091 & 0.072 & [0.045, 0.095] \\
Transfer tenure (days) & 0.363 & 0.294 & [0.274, 0.316] \\
Active months & 0.377 & 0.327 & [0.307, 0.348] \\
\addlinespace
\multicolumn{4}{l}{\textit{AWP}} \\
Inbound counterparty ratio & -0.269 & -0.228 & [-0.250, -0.207] \\
Transfer tenure (days) & 0.685 & 0.507 & [0.493, 0.521] \\
Active months & 0.704 & 0.555 & [0.540, 0.569] \\

\bottomrule
\end{tabular}
}
\end{table}


표~\ref{tab:alignment-sybil-stability}를 보면, AWP는 활동 기간(첫 기록부터 마지막 기록까지 지난 날 수)에서 $0.507$ 대 $0.294$로, 활동한 달 수에서 $0.555$ 대 $0.327$로 엔도스랭크보다 앞서요. 기간 전체에 퍼진 활동에 상을 주는 송금 점수라면 그럴 것이라고 예상할 수 있어요. 가족 평균에 대한 짝지은 차이도 구간이 0을 벗어나요(표~\ref{tab:tau-diff}, 대비~iv). 들어오는 상대 비율(받은 송금 수에 견준 서로 다른 보낸 쪽 수)에서는 AWP가 뚜렷하게 음수이고($-0.228$), 엔도스랭크는 약하게 양수예요($0.072$). 이 가족은 각 점수가 시간에 걸친 송금 정보를 얼마나 담고 있는지 보여 줘요.

표~\ref{tab:alignment-summary}의 보관해 둔 거래 가족에서, AWP는 청산(손실이 너무 커져 거래가 강제로 끝나는 일)에서 $0.141$ 대 $0.093$으로, 거래 성공에서 $0.236$ 대 $0.096$으로 엔도스랭크보다 앞서요. 위험 반대(값이 클수록 손실이 적은 지표) 가족에서는 둘 다 음수로, 엔도스랭크는 $-0.055$, AWP는 $-0.127$이에요. 손실 없이 닫은 비율은 거래량이 늘어도 커지지 않는 비율인데, 여기서 두 계수는 $-0.031$과 $-0.006$이에요.

\dissertationSubheading[sec:er-awp-holdout]{기간 밖}

표~\ref{tab:endorserank-awp-holdout}는 두 홀드아웃 기간에서 두 점수 사이의 짝지은 차이를 모아 두었어요. 스펜더(spender, 허락을 받아 대신 쓰는 쪽) 코호트 전체에서는, 라벨(나중에 채점에 쓰는 정답 값)이 세는 화살표 종류로 만든 점수를 $a$로 삼아요. 두 번째 묶음의 부분 집합에서는 두 행 모두 엔도스랭크를 $a$로 삼아요.

% Spring holdout and post hoc June-August scores -- regenerate with:
%   2-Dissertation-Draft-Works/wallet-reputation-experiments/scripts/run_holdout_eval.py --cohort spenders
%   2-Dissertation-Draft-Works/wallet-reputation-experiments/scripts/run_fresh_posthoc.py
% Generated by export_latex_results.py -- do not edit by hand unless re-export fails
\begin{table}[htbp]
\centering
\caption[Paired bootstrap differences between EndorseRank and AWP out of window.]{Paired bootstrap differences between EndorseRank and AWP out of window. None was fixed in advance: the spring rows were computed after the spring labels were known, and EndorseRank was scored on the June--August cohorts after the registered evaluation (Table~\ref{tab:fresh-posthoc}). The second block is the subset of Table~\ref{tab:holdout-receiving}. Within each block $a$ and $b$ share the wallet resamples (400, seed 42).}
\label{tab:endorserank-awp-holdout}
\fitwidth{%
\begin{tabular}{llccccc}
\toprule
Contrast ($a$ minus $b$) & Label & $\tau_a$ & $\tau_b$ & $\Delta\tau$ & 95\% CI of $\Delta\tau$ & $P(\Delta\tau>0)$ \\
\midrule
\multicolumn{7}{l}{\emph{Spring holdout, spenders ($n=1{,}335$)}} \\
EndorseRank minus AWP & new approval pairs & 0.394 & 0.123 & +0.271 & [0.207, 0.332] & 1.000 \\
AWP minus EndorseRank & new transfer senders & 0.344 & 0.283 & +0.061 & [0.026, 0.098] & 1.000 \\
\midrule
\multicolumn{7}{l}{\emph{Spring holdout, spenders that had received a transfer before the freeze ($n=977$)}} \\
EndorseRank minus AWP & new approval pairs & 0.340 & 0.298 & +0.042 & [0.003, 0.085] & 0.978 \\
EndorseRank minus AWP & new transfer senders & 0.358 & 0.388 & -0.029 & [-0.070, 0.011] & 0.085 \\
\midrule
\multicolumn{7}{l}{\emph{June--August, spenders ($n=1{,}964$)}} \\
EndorseRank minus AWP & new approval pairs & 0.333 & 0.118 & +0.215 & [0.171, 0.270] & 1.000 \\
AWP minus EndorseRank & new transfer senders & 0.319 & 0.255 & +0.064 & [0.036, 0.088] & 1.000 \\
\midrule
\multicolumn{7}{l}{\emph{June--August, traders with at least three closes ($n=1{,}402$)}} \\
EndorseRank minus AWP & liquidation-free close rate & 0.081 & 0.045 & +0.035 & [-0.017, 0.087] & 0.925 \\
\bottomrule
\end{tabular}
}
\end{table}


봄 스펜더에서 엔도스랭크는 새 허락 쌍(새로 생긴 주인–스펜더 허락)을 $\tau=0.394$(구간 $[0.360,0.431]$)로 줄 세우고, AWP는 $0.123$이에요. 짝지은 차이는 $+0.271$($[0.207,0.332]$)이에요. 새 송금자(새로 송금을 보내온 주소)에서는 순서가 뒤집혀요. AWP는 $0.344$이고 엔도스랭크는 $0.283$이어서, 차이는 $+0.061$($[0.026,0.098]$)이에요. 모든 라벨은 점수를 매긴 모든 이벤트보다 나중에 일어나요. 그래서 이 코호트에서는 화살표 종류마다, 다른 쪽에는 없는 자기 구성개념에 대한 정보를 담고 있어요. 첫 번째 차이의 많은 부분은 다루는 범위로 설명돼요. AWP는 동결일 전에 송금 화살표가 없는 스펜더 318개에 0점을 주는데, 이 스펜더들은 나머지보다 새 허락 쌍을 더 자주 얻었어요. 송금을 받은 적이 있는 스펜더 977개에서는, 새 허락 쌍에서 엔도스랭크가 앞서는 폭이 $+0.042$($[0.003,0.085]$)로 줄어요. 그리고 새 송금자에서는 두 점수가 비슷해요($-0.029$, $[-0.070,0.011]$). 송금 정보를 전혀 쓰지 않는 받은 허락 수(허락 입차수)도 새 송금자에서 AWP보다 앞서요($0.404$ 대 $0.344$, 짝지은 검사는 하지 않음).

6월부터 8월까지의 기간에서도 같은 모습이 되풀이돼요. 이 기간에는 등록해 둔 평가를 마친 뒤에 엔도스랭크 점수를 매겼어요. 엔도스랭크는 새 허락 쌍을 $0.333$으로, AWP는 $0.118$로 줄 세워요($+0.215$, $[0.171,0.270]$). AWP는 새 송금자를 $0.319$로, 엔도스랭크는 $0.255$로 줄 세워요($+0.064$, $[0.036,0.088]$). 트레이더(거래하는 사람)의 청산 없이 닫은 비율에서는 엔도스랭크가 $0.081$, AWP가 $0.045$예요. 이 차이는 0과 구별할 수 없어요($+0.035$, $[-0.017,0.087]$).

탐색적(미리 정하지 않고 살펴본) 봄 실행의 매칭 트레이더(표~\ref{tab:holdout-traders})에서, AWP는 이익을 낸 닫기 횟수를 $0.210$(구간 $[0.181,0.238]$)으로 줄 세우고 엔도스랭크는 $0.072$($[0.041,0.111]$)예요. 실현한 이익은 AWP가 $0.191$($[0.159,0.219]$), 엔도스랭크가 $0.048$($[0.017,0.082]$)이에요. 이익을 낸 닫기의 비율에서는 어느 점수도 아무것도 보여 주지 못해요. AWP는 $0.016$($[-0.015,0.044]$)이고 엔도스랭크는 $-0.016$($[-0.053,0.024]$)이에요. AWP가 앞서는 것은 횟수와 합계인데, 이것들은 지갑이 거래를 많이 할수록 커져요. 그리고 같은 두 라벨에서 송금 입차수도 AWP와 거의 비슷하게 잘해요.

% !TEX root = ../main.tex
\section{주장과 증거 지도}
\label{ch:appendix-claims}

이 부록은 가설(H), 연구 질문(RQ), 그리고 자료로 뒷받침하는 주장(C)을 하나하나 \ref{ch:results}장의 관련 증거와 이어 줘요.

\dissertationSubheading[sec:map-h1]{H1 / RQ1: 순위 갈라짐}

\begin{itemize}
\item 내용: 같은 코호트(평가하는 지갑 묶음)에서 엔도스랭크(EndorseRank)와 AWP(적응형 가중 페이지랭크)는 지갑을 서로 다르게 줄 세워요.
\item 증거: \ref{sec:rank-divergence}절; 표~\ref{tab:tie-sensitivity}(마지막 행이 두 방법 사이의 $\tau_b$를 구간과 함께 보여 주고, 각 그래프 밖의 지갑을 외톨이 점(화살표가 하나도 없는 점)으로 남겨 둔 경우도 보여 줘요); 그림~\ref{fig:rank-scatter}와~\ref{fig:score-distributions}.
\item 해석: 두 순위는 같지 않고, 둘이 맞는 부분은 엔도스랭크가 0점을 주는 지갑들에서 나와요. 목표(O) O1의 기준은 채워졌어요(표~\ref{tab:objectives-achievement}).
\end{itemize}

\dissertationSubheading[sec:map-h2]{H2 / RQ2: 같은 기간 일치도}

\begin{itemize}
\item 내용: 같은 기간 안에서, 각 점수는 자기가 걷는 화살표로 만든 점검과 맞아요.
\item 증거: 모든 점수의 가족 평균(비슷한 지표 묶음에 대한 평균값)과 그 신뢰구간(믿을 만한 범위, 표~\ref{tab:alignment-hybrid}); 그래프 밖의 지갑을 외톨이 점으로 남겨 두었을 때의 같은 기간 계수(표~\ref{tab:tie-sensitivity}).
\item 참고: 대리 지표(대신 재는 값)가 그래프와 같은 이벤트에서 나오므로, 이 일치는 처음부터 어느 정도 생길 수밖에 없어요. 엔도스랭크의 허락 계수와 송금 계수는 천장(얻을 수 있는 가장 큰 값)이 서로 달라요(\ref{sec:same-window-alignment}절). 그래서 이 주장은 엔도스랭크가 둘 가운데 어느 쪽을 더 가깝게 따르는지에 대해서는 아무것도 말하지 않아요. 그리고 외톨이 점 규칙은 결과를 안 뒤에 시험해 본 것이에요. 이 비교들은 하나도 미리 정해 두지 않았어요. 엔도스랭크와 AWP 사이의 짝지은 차이와 대리 지표별 결과는 부록~\ref{ch:appendix-er-awp}에 있어요.
\end{itemize}

\dissertationSubheading[sec:map-h3]{H3 / RQ3: 계산 비용}

\begin{itemize}
\item 내용: $n$이 같을 때, 점수의 계산 시간은 그 그래프의 화살표 수를 따라가요. 그래서 가장 듬성듬성한 그래프를 쓰는 엔도스랭크가 가장 빨리 돌아가요.
\item 증거: 표~\ref{tab:benchmark-runtime}(실행 시간, 표준편차, 메모리, 반복 수, $|V|$, $|E|$); 표~\ref{tab:benchmark-scaling-incohort}와~\ref{tab:benchmark-scaling}; 그림~\ref{fig:runtime-scaling}.
\item 참고: 확장 풀 단계는 뽑은 주소를 더하지만, 매칭 코호트 그래프 너머로 화살표는 더하지 않아요. 그래서 풀이 프로그램 그래프의 점(노드)도 더하지 않아요(\ref{sub:runtime-scaling}절). 이 주장은 화살표 수에 대한 것으로만 좁혀져요.
\end{itemize}

\dissertationSubheading[sec:map-rq4]{RQ4: 시간 홀드아웃}

\begin{itemize}
\item 내용: 점수를 $t_1$에 고정하면, 그 점수들은 스펜더(spender, 허락을 받아 대신 쓰는 쪽) 코호트($n=1{,}335$)에서 $t_1$ 뒤부터 $t_2$까지 나타나는 새 주인--스펜더 허락과 새 송금자를 예측해요(미리 맞혀요). 점수가 어느 라벨(나중에 채점에 쓰는 정답 값)을 잘 줄 세우는지는 화살표 종류가 정해요. 하지만 어떤 페이지랭크도 자기 화살표 종류의 $t_1$ 원시 차수(화살표를 그냥 센 수)보다 어느 라벨이든 더 잘 예측하지는 못해요.
\item 증거: 표~\ref{tab:holdout-spenders}와~\ref{tab:holdout-diff}(\ref{sec:holdout-results}절); 6월부터 8월까지의 기간은 표~\ref{tab:fresh-posthoc}; 동결일(점수를 고정한 날) 전에 송금을 받은 적이 있는 스펜더는 표~\ref{tab:holdout-receiving}.
\item 단서: 두 라벨은 나중 기간에 관찰한 구성개념(점수가 재려고 하는 생각) 그 자체예요(새 허락, 새 송금자). 신용이나 거래 결과가 아니에요. 그리고 둘 다 새로 들어온 것을 센 수예요. RQ4는 연구 계획서에 있던 거래 성공 검증을 대신했어요(\ref{sub:gf-pr}절). C-PR($\lambda=1$)의 증가분(자기 차수보다 나아진 몫)을 빼면, 이 비교들은 하나도 미리 정해 두지 않았어요. 엔도스랭크와 AWP 사이의 차이는 부록~\ref{ch:appendix-er-awp}에 있어요.
\end{itemize}

\dissertationSubheading[sec:map-rq5]{RQ5: 결합 연산자}

\begin{itemize}
\item 내용: 9월 14일의 규칙에 따르면, 금액을 그대로 쓴 허락 층과 송금 층을 점마다 볼록 결합(정해진 비율로 더해 섞기)한 것(C-PR)은 한 층만 걷는 두 걷기를 둘 다 이기지는 못해요. 그리고 씨앗을 쓴 빼 보기 실험(한 부분만 바꿔 그 효과를 보는 실험), 곧 S-PR은 송금 걷기를 그대로 다시 만들어 내요. 봄 홀드아웃(나중 기간을 떼어 두고 맞혀 보는 시험)의 탐색적(미리 정하지 않고 살펴본) 비교에서는 C-PR이 송금 걷기만 한 것보다 나았지만, 등록 재현(계획을 먼저 공개해 두고 새 자료로 다시 해 보기)은 이것을 확인해 주지 않아요. F1과 F3은 통과하고, F2는 통과하지 못해요. AWP를 비교 대상으로 삼으면 F2는 더 큰 차이로 통과하지 못해요.
\item 증거, 9월 14일의 규칙: 주된 부분은 표~\ref{tab:holdout-diff}, 두 번째 부분은 표~\ref{tab:alignment-hybrid}와~\ref{tab:tau-diff-hybrid}(\ref{sec:coupled-results}절).
\item 증거, 송금 걷기만 한 것과의 비교: 표~\ref{tab:holdout-diff}의 탐색적 봄 대비, 그리고 표~\ref{tab:fresh-contrasts}의 등록한 판정(F1--F3; 그 옆에 F4--F6도 함께 보고; \ref{sub:fresh-results}절).
\item 증거, AWP와의 비교: 표~\ref{tab:fresh-contrasts}의 등록한 민감도 분석(조건을 바꿔 결과가 얼마나 움직이는지 보기), 그리고 표~\ref{tab:holdout-diff}에 있는, 라벨을 본 뒤에 계산한 봄 대비.
\item 증거, 비용: 표~\ref{tab:benchmark-runtime}의 실행 시간 행.
\item 단서: 고정한 $\lambda\in\{0.25,0.5,0.75\}$, 금액을 그대로 쓴 층, 고르게 순간 이동하기, 이동(전이) 단계에서 섞기, 그리고 체인 하나로만 시험했어요. 9월 14일의 규칙과 $\lambda$ 값들은 결합 점수를 하나라도 계산하기 전에 설정 파일에 정해 두었어요(\ref{sec:holdout-protocol}절). 송금 걷기만 한 것과 비교하는 규칙은 2026년 9월 28일에 6월부터 8월까지의 기간을 위해 썼어요. 그리고 추출하기 전인 2026년 10월 1일에 민감도 분석 두 개와 함께 등록했는데, 그중 하나는 AWP를 비교 대상으로 삼아요(\ref{sub:fresh-holdout}절). 등록 문서에서는 한 층만 걷는 두 걷기를 엔도스랭크와 AWP라고 불러요.
\end{itemize}

\dissertationSubheading[sec:map-c1c4]{주장 C1--C5}

\begin{itemize}
\item C1: H3/RQ3 아래에 적은 실행 시간 증거예요. 시간 분석 실행(어느 부분에 시간이 드는지 재 본 실행)과 화살표 수 비율은 \ref{sec:benchmark-results}절과 부록~\ref{sec:pr-sparsity}에 있어요. 이 주장은 엔도스랭크에만 맞는 말이에요. C-PR은 적어도 AWP만큼 비용이 들기 때문이에요.
\item C2: 모든 점수가 같은 기간에 자기 화살표의 점검과 맞는 것이에요(표~\ref{tab:alignment-hybrid}). 이 일치는 처음부터 어느 정도 생길 수밖에 없고, 엔도스랭크의 경우에는 그래프 밖 지갑에 대한 규칙이 정해요(표~\ref{tab:tie-sensitivity}).
\item C3: H1/RQ1 아래에 적은 순위 갈라짐이에요(표~\ref{tab:tie-sensitivity}, 그림~\ref{fig:rank-scatter}).
\item C4: 스펜더 홀드아웃에서 모든 페이지랭크를 자기 화살표 종류의 원시 차수와 비교한 것이에요(표~\ref{tab:holdout-spenders}와~\ref{tab:holdout-diff}). 여기에 다루는 범위 점검(표~\ref{tab:holdout-receiving}), 트레이더(거래하는 사람)에 대한 한계를 보여 주는 탐색적 트레이더 실행(표~\ref{tab:holdout-traders}), 그리고 이 결과를 되풀이한 6월부터 8월까지의 기간(표~\ref{tab:fresh-posthoc})이 함께 있어요. 강건성 점검(조건을 바꿔도 결과가 버티는지 보기)과 표본 정의 민감도는 표~\ref{tab:robustness-damping}--\ref{tab:robustness-sample-size}, 표~\ref{tab:endorserank-restarts}, 그림~\ref{fig:damping-sensitivity}에 보고했어요.
\item C5: 결합은 한 층만 걷는 두 걷기를 이기지 못해요. 이것은 조건이 정해진 부정적인 결과예요. 그리고 등록 재현은 송금 걷기만 한 것보다 낫다는 탐색적 이득을 확인해 주지 않아요(표~\ref{tab:fresh-contrasts}). 위의 RQ5와 \ref{sec:coupled-results}절을 보세요.
\end{itemize}

\dissertationSubheading[sec:map-nonclaims]{주장하지 않는 것}

이 연구는 다음을 주장하지 \emph{않아요}.

\begin{itemize}
\item 어느 점수가 평판을 더 잘 잰다거나, AWP가 이제 쓸모없어졌다는 것(부록~\ref{ch:appendix-er-awp}의 엔도스랭크와 AWP 비교는 구성개념을 따라가고, 다루는 범위에 따라 달라져요);
\item 엔도스랭크가 송금 가족보다 허락 가족과 더 가깝게 맞는다는 것(두 계수는 천장이 서로 달라요);
\item 이 논문의 어떤 페이지랭크가 자기 구성개념이 앞으로 자라는 것을 $t_1$의 원시 차수보다 더 잘 미리 맞힌다는 것;
\item 값 매개변수 $b=1$이나 엔도스랭크의 고르게 다시 출발하기가 가장 좋은 선택이라는 것(둘 다 이 연구가 정했고, 재시작은 이 자료에서 두 규칙을 비교한 뒤에 정했어요);
\item C-PR이 송금만 걷는 것보다 낫다는 것(등록한 규칙이 새 송금자에서 통과하지 못했어요), 또는 AWP보다 낫다는 것(두 기간 모두 새 송금자에서 AWP에 뒤져요);
\item 두 화살표 종류를 어떻게 섞어도 한 층 방법을 이길 수 없다는 것. 여기서 시험한, $\lambda$를 고정한 점별 볼록 결합이 이기지 못한다는 것만 말해요;
\item 화살표 수가 코호트 그래프보다 많아질 때 실행 시간이 유리하게 늘어난다는 것;
\item 다시 해 보지 않고도 결과가 아비트럼 원(Arbitrum One)과 연구 기간 밖으로 넓혀진다는 것;
\item 온체인(블록체인 위의) 점수가 담보나 법적 신분 확인 제도를 대신한다는 것.
\end{itemize}

\dissertationSubheading[sec:map-artifacts]{재현용 결과물}

\begin{enumerate}
\item 엔도스랭크, AWP, C-PR($\lambda\in\{0,0.25,0.5,0.75,1\}$), S-PR의 점수를 하나로 합친 지갑 순위 표.
\item 스펜더 코호트와 매칭 트레이더에 대한, 컴퓨터가 읽을 수 있는 홀드아웃 요약. 짝지은 대비와 주된 기준의 판정이 들어 있어요(\texttt{data/3-published-results-for-thesis/spring-holdout-2026-03-to-2026-05/} 아래에 커밋됨).
\item 컴퓨터가 읽을 수 있는 평가 요약: 기준이 되는 같은 기간 숫자, 실행마다 잰 시간, 부트스트랩(다시 뽑기) 절차(\texttt{data/3-published-results-for-thesis/} 아래에 커밋됨).
\item 등록 재현: 2026년 10월 1일에 커밋한 그대로의 등록 파일과 계획서, 두 파일의 해시(파일 내용으로 만든 지문 같은 값)가 적힌 추출 목록 파일(매니페스트), 모든 계수·대비·판정이 담긴 평가 요약, 그리고 같은 코호트에서 사후에(라벨을 본 뒤에) 매긴 엔도스랭크 점수(\texttt{data/3-published-results-for-thesis/} 안의 \texttt{config-copies/}와 \texttt{registered-replication-2026-06-to-2026-08/} 아래에 커밋됨).
\item 내보낸 표 조각(\texttt{results/tables/}).
\item 내보낸 벡터 그림(크게 키워도 깨지지 않는 그림, \texttt{Figures/generated/}).
\item 꺼낸 로그와 그것으로 만든 표(\texttt{data/1-raw-blockchain-logs/}와 \texttt{data/2-processed-tables-and-evaluations/}). 100\,MiB가 넘는 파일은 빠져 있지만, 한 달 단위 파일로 다시 만들 수 있어요.
\item 보관해 둔 확장 기준선(비교 기준) 행렬(\texttt{data/4-archived-extended-baselines/}, 선택 사항).
\end{enumerate}

표와 그림은 이렇게 다시 만들어요. 실제 자료로 강건성 점검 실행까지 포함해 평가 처리 과정을 돌리고, 표를 내보낸 다음, 같은 결과물로 그림을 다시 그려요(부록~\ref{sec:pipeline-stages}).

% !TEX root = ../main.tex
\section{윤리 승인 증명서}
\label{ch:appendix-ethics}

\dissertationSubheading[sec:ethics-certificate]{증명서}

이 연구는 공개된 블록체인 자료만 쓰고, 사람 참여자는 없어요(\ref{sec:ethics}절). Unisa(남아프리카 대학교) 컴퓨팅 학부 / 농업·환경과학 대학의 연구 윤리 심의 위원회가 내준 윤리 승인(연구가 윤리 기준을 지킨다는 허가) 증명서를 아래에 옮겨 실어요.

% TODO (user action): replace the placeholder box with the certificate PDF once
% the reference number, date and scanned document are available, e.g.
%   \includegraphics[width=\textwidth,page=1]{Figures/ethics-clearance.pdf}
% Do not invent a reference number or date.
\begin{center}
\fbox{\parbox{0.9\textwidth}{\centering\vspace{2em}
윤리 승인 증명서\\[0.5em]
참조 번호: [Unisa SoC/CAES 윤리 승인 번호: 미정]\\
발급일: [미정]\\[0.5em]
(스캔한 증명서를 여기에 넣을 예정이에요.)
\vspace{2em}}}
\end{center}

